% GENERATED FILE. Edit canonical lessons, then run npm run generate:books.
% canonical-source-hash: fnv1a64:b2adfb65b74d7107
% canonical-lessons: TE-C377-vipattu, TE-C377-alajadi, TE-C377-raksana, TE-C377-gasti, TE-C377-carya

\chapter{Disaster}
\label{ch:te-disaster}
\hlchaptermodality{377}

\begin{chapteropening}
\textbf{By the end of this chapter:} \emph{I can say a disaster, a commotion, protection, a patrol and a measure.}

Complete the last lesson of chapter 377: I can say a disaster, a commotion, protection, a patrol and a measure.
\end{chapteropening}

\section[\texorpdfstring{\te{విపత్తు} (vipattu)}{విపత్తు (vipattu)}]{\te{విపత్తు} (vipattu) — a disaster; a calamity}

\label{lesson:TE-C377-vipattu}



\begin{quote}
Before the new one: say the Telugu for an episode, then the Telugu for by chance.
\end{quote}

\subsection*{You'll want to know: \te{విపత్తు}}
\textbf{\te{విపత్తు}} — \emph{vipattu} — \textquotedblleft{}a disaster; a calamity\textquotedblright{}.

\begin{quote}
\te{ఆ} \te{విపత్తులో} \te{చాలా} \te{ఇళ్ళు} \te{పోయాయి}.
\end{quote}

\begin{quote}
\emph{ā vipattulō cālā iḷḷu pōyāyi.}
\end{quote}

\begin{quote}
Many houses were lost in that disaster.
\end{quote}

Plural: \te{విపత్తులు}, \emph{vipattulu}.

\te{ప్రకృతి} \te{విపత్తు}, a natural disaster.

\begin{cousinweb}
\textbf{\te{విపత్తు}} (\emph{vipattu}) is borrowed from Sanskrit \emph{vipatti}, \textquotedblleft{}misfortune, failure\textquotedblright{} (CDIAL 11788).
\end{cousinweb}

\begin{grammarlens}[title={What you've built}]
One more word for telling what happened.
\end{grammarlens}

\subsection*{Your turn}
\begin{itemize}
\raggedright
  \item \emph{Say it:} \emph{vipattu}
  \item \emph{Say it:} \emph{vipattu}, once more
  \item \emph{Say it:} \emph{ā vipattulō cālā iḷḷu pōyāyi.} — Many houses were lost in that disaster.
  \item \emph{Recall:} say the Telugu for an episode, then the Telugu for by chance, then say \emph{vipattu} again
\end{itemize}

\begin{tcolorbox}[breakable,colback=teal!4,colframe=teal!35!black,title={Before you move on}]
What does \te{విపత్తు} mean? (\textquotedblleft{}A disaster; a calamity\textquotedblright{}.) Say it once more.
\end{tcolorbox}
\section[\texorpdfstring{\te{అలజడి} (alajaḍi)}{అలజడి (alajaḍi)}]{\te{అలజడి} (alajaḍi) — a commotion; unrest}

\label{lesson:TE-C377-alajadi}



\begin{quote}
Before the new one: say the Telugu for by chance, then the Telugu for a disaster.
\end{quote}

\subsection*{You'll want to know: \te{అలజడి}}
\textbf{\te{అలజడి}} — \emph{alajaḍi} — \textquotedblleft{}a commotion; unrest\textquotedblright{}.

\begin{quote}
\te{రాత్రి} \te{వీధిలో} \te{అలజడి} \te{మొదలైంది}.
\end{quote}

\begin{quote}
\emph{rātri vīdhilō alajaḍi modalaindi.}
\end{quote}

\begin{quote}
A commotion started in the street at night.
\end{quote}

\begin{cousinweb}
\textbf{\te{అలజడి}} (\emph{alajaḍi}) is native Dravidian: DEDR 236 lists Telugu \emph{alajaḍi}, \textquotedblleft{}sorrow, affliction, calamity\textquotedblright{}, under the word for distress; today it names a stir of unrest.
\end{cousinweb}

\begin{grammarlens}[title={What you've built}]
One more word for telling what happened.
\end{grammarlens}

\subsection*{Your turn}
\begin{itemize}
\raggedright
  \item \emph{Say it:} \emph{alajaḍi}
  \item \emph{Say it:} \emph{alajaḍi}, once more
  \item \emph{Say it:} \emph{rātri vīdhilō alajaḍi modalaindi.} — A commotion started in the street at night.
  \item \emph{Recall:} say the Telugu for by chance, then the Telugu for a disaster, then say \emph{alajaḍi} again
\end{itemize}

\begin{tcolorbox}[breakable,colback=teal!4,colframe=teal!35!black,title={Before you move on}]
What does \te{అలజడి} mean? (\textquotedblleft{}A commotion; unrest\textquotedblright{}.) Say it once more.
\end{tcolorbox}
\section[\texorpdfstring{\te{రక్షణ} (rakṣaṇa)}{రక్షణ (rakṣaṇa)}]{\te{రక్షణ} (rakṣaṇa) — protection; safety}

\label{lesson:TE-C377-raksana}



\begin{quote}
Before the new one: say the Telugu for a disaster, then the Telugu for a commotion.
\end{quote}

\subsection*{You'll want to know: \te{రక్షణ}}
\textbf{\te{రక్షణ}} — \emph{rakṣaṇa} — \textquotedblleft{}protection; safety\textquotedblright{}.

\begin{quote}
\te{ప్రభుత్వం} \te{ప్రజల} \te{రక్షణ} \te{కోసం} \te{పోలీసులను} \te{పంపింది}.
\end{quote}

\begin{quote}
\emph{prabhutvaṁ prajala rakṣaṇa kōsaṁ pōlīsulanu pampindi.}
\end{quote}

\begin{quote}
The government sent police for the people's protection.
\end{quote}

Also defence, as in the name of a ministry.

\begin{cousinweb}
\textbf{\te{రక్షణ}} (\emph{rakṣaṇa}) is borrowed from Sanskrit \emph{rakṣaṇa}, \textquotedblleft{}guarding, protecting\textquotedblright{} (CDIAL 10546).
\end{cousinweb}

\begin{grammarlens}[title={What you've built}]
One more word for telling what happened.
\end{grammarlens}

\subsection*{Your turn}
\begin{itemize}
\raggedright
  \item \emph{Say it:} \emph{rakṣaṇa}
  \item \emph{Say it:} \emph{rakṣaṇa}, once more
  \item \emph{Say it:} \emph{prabhutvaṁ prajala rakṣaṇa kōsaṁ pōlīsulanu pampindi.} — The government sent police for the people's protection.
  \item \emph{Recall:} say the Telugu for a disaster, then the Telugu for a commotion, then say \emph{rakṣaṇa} again
\end{itemize}

\begin{tcolorbox}[breakable,colback=teal!4,colframe=teal!35!black,title={Before you move on}]
What does \te{రక్షణ} mean? (\textquotedblleft{}Protection; safety\textquotedblright{}.) Say it once more.
\end{tcolorbox}
\section[\texorpdfstring{\te{గస్తీ} (gastī)}{గస్తీ (gastī)}]{\te{గస్తీ} (gastī) — a patrol}

\label{lesson:TE-C377-gasti}



\begin{quote}
Before the new one: say the Telugu for a commotion, then the Telugu for protection.
\end{quote}

\subsection*{You'll want to know: \te{గస్తీ}}
\textbf{\te{గస్తీ}} — \emph{gastī} — \textquotedblleft{}a patrol\textquotedblright{}.

\begin{quote}
\te{రాత్రంతా} \te{పోలీసులు} \te{గస్తీ} \te{తిరిగారు}.
\end{quote}

\begin{quote}
\emph{rātrantā pōlīsulu gastī tirigāru.}
\end{quote}

\begin{quote}
The police patrolled all night.
\end{quote}

\te{గస్తీ} \te{తిరుగు}, to go on patrol.

\begin{cousinweb}
\textbf{\te{గస్తీ}} (\emph{gastī}) is a loan from Persian \emph{gašt}, \textquotedblleft{}a walking round, a patrol\textquotedblright{}.
\end{cousinweb}

\begin{grammarlens}[title={What you've built}]
One more word for telling what happened.
\end{grammarlens}

\subsection*{Your turn}
\begin{itemize}
\raggedright
  \item \emph{Say it:} \emph{gastī}
  \item \emph{Say it:} \emph{gastī}, once more
  \item \emph{Say it:} \emph{rātrantā pōlīsulu gastī tirigāru.} — The police patrolled all night.
  \item \emph{Recall:} say the Telugu for a commotion, then the Telugu for protection, then say \emph{gastī} again
\end{itemize}

\begin{tcolorbox}[breakable,colback=teal!4,colframe=teal!35!black,title={Before you move on}]
What does \te{గస్తీ} mean? (\textquotedblleft{}A patrol\textquotedblright{}.) Say it once more.
\end{tcolorbox}
\section[\texorpdfstring{\te{చర్య} (carya)}{చర్య (carya)}]{\te{చర్య} (carya) — an action; a measure}

\label{lesson:TE-C377-carya}



\begin{quote}
Before the new one: say the Telugu for protection, then the Telugu for a patrol.
\end{quote}

\subsection*{You'll want to know: \te{చర్య}}
\textbf{\te{చర్య}} — \emph{carya} — \textquotedblleft{}an action; a measure\textquotedblright{}.

\begin{quote}
\te{ప్రభుత్వం} \te{వెంటనే} \te{చర్యలు} \te{తీసుకుంది}.
\end{quote}

\begin{quote}
\emph{prabhutvaṁ veṇṭanē caryalu tīsukundi.}
\end{quote}

\begin{quote}
The government took action at once.
\end{quote}

Plural: \te{చర్యలు}, \emph{caryalu}.

\te{చర్య} \te{తీసుకో}, to take action.

\begin{cousinweb}
\textbf{\te{చర్య}} (\emph{carya}) is borrowed from Sanskrit \emph{caryā}, \textquotedblleft{}conduct, a going about\textquotedblright{} (CDIAL 4710).
\end{cousinweb}

\begin{grammarlens}[title={What you've built}]
One more word for telling what happened.
\end{grammarlens}

\subsection*{Your turn}
\begin{itemize}
\raggedright
  \item \emph{Say it:} \emph{carya}
  \item \emph{Say it:} \emph{carya}, once more
  \item \emph{Say it:} \emph{prabhutvaṁ veṇṭanē caryalu tīsukundi.} — The government took action at once.
  \item \emph{Recall:} say the Telugu for protection, then the Telugu for a patrol, then say \emph{carya} again
\end{itemize}

\begin{tcolorbox}[breakable,colback=teal!4,colframe=teal!35!black,title={Before you move on}]
What does \te{చర్య} mean? (\textquotedblleft{}An action; a measure\textquotedblright{}.) Say it once more.
\end{tcolorbox}
